\documentclass[12pt]{article}
\usepackage[utf8]{inputenc}
\usepackage[T1]{fontenc}
\usepackage{hyperref}
\title{Swahili Language Resources for K-12 Curriculum Development}
\author{Compiled by Claude Sonnet 5.5 (claude-sonnet-5-5) for Andr\'as Kornai}
\date{}
\begin{document}
\maketitle

\section{Introduction}
ISO 639-3 \texttt{swa} is the Swahili macrolanguage code. Standard Kiswahili is used as a medium of instruction in Tanzanian primary schools, while English is the medium in secondary schools. The searches I ran did not return a current catalogue of official textbooks, so this report lists only items that appeared in search results. Related codes such as \texttt{swc} (Congo Swahili) and \texttt{swb} (Maore Comorian) are covered separately and are not mixed in here.

\section{Primary Educational Resources}
\begin{itemize}
\item Zanzibar Kiswahili syllabus for lower primary (Standards I--III), Zanzibar Institute of Education: \url{https://moez.go.tz/dox/syllabus/kiswahili_lower_primary.pdf}. The document carries a copyright notice restricting reproduction.
\item Kiswahili syllabus S1--S6, Zanzibar, 2009: \url{https://maktaba.tetea.org/study-aids/primary/Syllabus - Kiswahili - S1-S6 Zanzibar - 2009 (Kiswahili).pdf}
\item USAID-hosted document from the search results, described as a Standard One Kiswahili reading teaching guide: \url{https://pdf.usaid.gov/pdf_docs/PA00KH6S.pdf}
\item Columbia University Libraries guide to Swahili resources, which lists the African Storybook among them: \url{https://library.columbia.edu/libraries/global/virtual-libraries/african_studies/languages/swahili.html}
\item Evaluation of A-level Kiswahili textbooks under Tanzania's 1991 textbook policy (University of Dar es Salaam): \url{https://libraryrepository.udsm.ac.tz/items/599be783-4311-40f4-a597-09b826ecef6e}
\item Article in \textit{Kioo cha Lugha} on Kiswahili textbooks written by the Tanzania Institute of Education (secondary): \url{https://journals.udsm.ac.tz/index.php/kcl/article/view/6138/4969}
\end{itemize}

\section{Linguistic and Reference Materials}
\begin{itemize}
\item Lexikos article on Swahili dictionaries, which discusses the TUKI dictionaries, \textit{Kamusi ya Kiswahili Sanifu} and online dictionaries: \url{https://lexikos.journals.ac.za/pub/article/view/1798/1036}
\item Second Lexikos article from the same search: \url{https://lexikos.journals.ac.za/pub/article/view/1273/786}
\item Nordic Journal of African Studies review that mentions the Helsinki Corpus of Swahili: \url{https://njas.fi/njas/article/download/679/502/1243}
\item Hunspell Swahili spell-checker README: \url{https://git.stg.centos.org/source-git/hunspell-sw/blob/c8/f/README_sw_TZ.txt}
\end{itemize}

\section{Religious and Community Resources}
\begin{itemize}
\item ScriptureEarth index of Swahili scriptures, including Bible.is text and audio: \url{https://scriptureearth.org/00i-Scripture_Index.php?iso=swh}
\item YouVersion Swahili Bible versions: \url{https://bible.com/languages/swh}
\item Swahili Christian resources page listing an audio Bible: \url{https://swahilichristian.missionresources.com/audio-bible.html}
\end{itemize}

\section{Digital and Audio Resources}
\begin{itemize}
\item Audio scripture on Bible.is, reached through the ScriptureEarth page above.
\item African Storybook, listed on the Columbia guide above (the guide's page is the only URL I saw for it).
\end{itemize}

\section{Curriculum Development Context}
Kiswahili is the language of instruction in Tanzanian primary schools and is taught as a subject in secondary schools. The sources found are mostly syllabi, academic evaluations and scripture sites. Several studies report that textbooks do not fully match the syllabus and that primary schools lack adequate Kiswahili texts. I did not find the Tanzania Institute of Education's current textbook catalogue in these results, and a curriculum project would need to check it directly.

\end{document}